\clearpage
\paragraph{Covariance and injection follow-up (v5.1.1).}
The new study tests two links in the approximation: whether the response
matrix describes the dependence between fitted masses, and whether a known
signal produces the expected local recovery and neighboring response.
The observed spectra, released limits and probabilities above are unchanged.
These are conditional diagnostics of the fixed analysis and generating
spectra, not a new global calibration.

For the extended-grid covariance check, the same first 256 archived
background-only Poisson spectra are evaluated at every coordinate and
centered by their empirical means. Their covariance
and correlation are compared with $D^\top D$ and $R$ from
Eqs.~\ref{v5-eq:field-response}--\ref{eq:v501-field-correlation}.
Bootstrap resampling keeps each whole scan together, preserving the
correlations among masses. The resulting uncertainties quantify Monte Carlo
precision; they do not include uncertainty in the stress-background model.
The eigenmodes describe correlated patterns in the scan. An effective
mode count is a covariance summary, not a calibrated number of search trials.
The RMS difference between empirical and response-predicted correlations
is 0.0556 for 2015, 0.0626 for 2016 and 0.0594 for 2021. The denser
2016 grid resolves the shape of the stress response, whose signed root
still spans $-17.86$ to $+13.71$; grid refinement does not repair that
background offset.

The 2015 pilot covers 19--100~MeV at integer masses. Above 90~MeV the
kernel coordinates are held at their 90~MeV values; each added mass still
has its own recomputed fit and source-bin response. The 2016 grid retains
39--180~MeV and adds half-MeV points from 70.5 to 80.5~MeV, with log-linear
interpolation of the saved kernel coordinates. The released 2021 sample
retains 50--250~MeV at integer masses. The grids therefore contain
82, 153 and 201 points. Re-evaluating the same 256 complete spectra at
all coordinates provides the new covariance blocks; unrelated local toys
are not joined. The earlier 1,000-scan native-grid benchmark is retained
separately. A timing-based amendment fixed this reduced covariance sample
and grid before inspecting their results. The extension remains a pilot
and does not enlarge the released observed search.

\fig{0.99}{../figures/v511_covariance_comparison.pdf}{Response-predicted
correlations and their empirical comparison using the same 256 complete
background-only Poisson spectra at all coordinates. Each row uses its
actual mass grid; difference panels have separate color scales. The
original 1,000-scan native-grid result is retained separately and is not
spliced into these extended matrices. This checks central dependence
under the specified stress reference, not the rare global tail.}
{fig:v511-covariance}
\FloatBarrier

\paragraph{Fixed-mass recovery and full-scan signal responses.}
The injection masses are inherited diagnostic anchors: 51, 78 and 95~MeV
for 2015; 42, 66, 76 and 92~MeV for 2016; and 71 and 92~MeV for 2021.
They cover the earlier 2016 stress excursions, the neighboring structures
illustrated by the 2021 65/71/78~MeV examples, representative peaks, and
the proposed 2015 extension. They were chosen to investigate those known
features, so they do not constitute an independent discovery selection.
Each mass is tested against both the archived coherent stress spectrum
and a full-support local GP control fitted to the observed sidebands
with that signal window excluded. The control depends on the injection
mass, so these local controls do not form a single global null. The same frozen reference
error defines the physical yields $A/\sigma_{\rm ref}=0,2,5$ for both
truths. Within a mass and truth, 64 Poisson background spectra are paired
across strengths; independent Poisson signal increments of two and three
reference errors build the size-two and size-five samples. Thus the
three strengths are correlated experiments, not independent ensembles.
Figure~\ref{fig:v511-injections} shows three representative anchors.
At 2016 76~MeV, the size-five injection gives mean signed root $-9.80$
under the stress truth and $+4.52$ under the local GP control; every
stress fit remains negative. At 2021 71~MeV the corresponding means are
$+3.14$ and $+4.37$. Adding a signal moves the distributions, but it does
not remove the background-dependent offset. With 64 experiments per
curve, these are small conditional recovery samples.

A deterministic injected spectrum is also scanned over every tested mass.
For a count perturbation $\delta n_i$ from the signal, the first-order
prediction from the background response is
\begin{equation}
 \Delta r_m^{\rm lin}=\sum_{i:B_i>0}D_{im}
                         \frac{\delta n_i}{\sqrt{B_i}}.
 \label{eq:v511-injection-projection}
\end{equation}
For the 2016 76~MeV size-five injection, the exact change is $+4.843$
at the injected mass and reaches $-3.589$ at 85~MeV. The largest difference
from the linear projection is 0.00924 in signed root over the scan.
Comparison with the exact difference of the injected and background-only
scans tests linearity and reveals neighboring peaks or deficits induced by
the shared interpolation. An off-mass dip can therefore accompany a
positive injection. This diagnostic does not establish the origin of a
feature in the observed spectrum. A covariance-matrix column describes
shared noise; it is not the signal-template projection in
Eq.~\ref{eq:v511-injection-projection}.
\fig{0.99}{../figures/v511_signed_injection_echoes.pdf}{Deterministic signal
injections into the coherent stress spectrum. Left: signed roots before
and after injection. Right: the size-five change with full retraining,
its first-order template projection, and a control holding the
background-only GP prediction and covariance fixed. The likelihood's
background nuisance parameters remain profiled in the last control.
Vertical lines mark the injected mass. Neighboring negative responses
are retained. These full scans contain no Poisson fluctuation.}
{fig:v511-echo}
\FloatBarrier

\fig{0.99}{../figures/v511_injection_distributions.pdf}{Fixed-mass Poisson
injection examples under local GP controls (top) and the archived stress
truth (bottom). Each curve uses 64 experiments, paired across the three
strengths within a mass and truth. Dotted lines show fits to unfluctuated
expectations. The horizontal axis is the signed fitted root; the labels
$A/\sigma_{\rm ref}=0,2,5$ specify injected yields, not achieved
significances. The same reference uncertainty is used for both truths
at a fixed mass. Persistent stress offsets are especially clear at
2016 76~MeV. No global discovery efficiency is assigned.}
{fig:v511-injections}
\FloatBarrier

The fixed-mass Poisson experiments measure conditional signal recovery
at their declared anchors. They do not scan each fluctuated injection
across the full mass range, and no new global false-positive threshold or
discovery power is inferred. The significance GP remains useful for
propagating a specified background response across the mass scan;
background qualification and global-tail validation remain separate
requirements.
\FloatBarrier
